% Copyright 2026 Sreeram Anil
% SPDX-License-Identifier: CC-BY-4.0
\chapter{Recommendations for an electric-aircraft BMS}
\label{ch:res-recommendations}

\section{The question the benchmark can answer}

The benchmark cannot say which estimator a particular aircraft should fly; that depends on the
cell, the mission, the sensors, the processor and the certification basis, and on measurements
that this report does not have. It can say, from controlled experiments on one representative
problem, which estimators are hurt by which faults and by how much, what each costs, and where
the accuracy of the whole system is set by something other than the estimator. This chapter
turns those findings into recommendations, stated as design rules with the evidence that
supports each, and closes with a reference architecture.

\section{Design rules}

\paragraph{Rule 1: calibrate the current sensor; no estimator will do it for you.}
A current-sensor offset and gain error are unobservable from one voltage channel
(\cref{ch:res-cell_ecm}, \code{current\_bias}), they are common-mode across a module
(\cref{ch:res-pack}), and a scale-factor error becomes a permanent capacity error of the same
size (\cref{ch:res-soh}). The estimators that partially compensate an offset --- MHE, the
disturbance observer, RBPF, AWTLS --- do so at a cost and only partially, and none compensates
a gain error. The nominal \SI{20}{\milli\ampere} and \SI{0.5}{\percent} of the benchmark are
already the floor of every estimator on the nominal mission; a flight BMS should budget its
current-sensor calibration against the SOC and SOH accuracy it needs, and re-anchor on an
OCV reading at rest whenever the mission allows.

\paragraph{Rule 2: keep the covariance honest.} The single most common failure in the
benchmark --- the EKF's \SI{12}{\percent} on \code{combined}, the sigma-point filters'
\code{sensor\_dropout}, the whole Kalman cluster on the electrochemical plant --- is a
covariance that collapses on the strength of a model that does not deserve the confidence. The
remedies are cheap and additive: a fading-memory factor sized from a memory length in samples
(\cref{ch:fadingekf}), a lower bound on the SOC variance, innovation-based adaptation of $\mR$
(\cref{ch:iaeakf,ch:sagehusaakf}) with a floor and a ceiling, and the Joseph-form update. An EKF
with the $\mR$ adaptation is the IAE-AKF, second of the ranking on the equivalent-circuit plant;
an EKF with the bounded memory is the Fading-EKF, first on the electrochemical plant; each
costs \SI{0.5}{\micro\second}. The two mechanisms address different failures --- the first a
sensor that is noisier than declared, the second a model that is worse than declared --- and
the per-cell estimator of the reference architecture below combines them; the combination was
not benchmarked as a separate estimator, and its tuning (the memory length in particular)
should be settled on the target cell.

\paragraph{Rule 3: use a capacity state, gate it on excitation, and let it run across flights.}
Only the joint and dual filters remove a capacity error (\code{aged\_cell}: \num{0.05} against
\num{1.5}--\num{4}\,\%), and only they track the fade over consecutive missions to the accuracy
of the current sensor (\cref{ch:res-soh}). The same filters are the worst of the field on the
ground-charge profile (\num{3.6}--\num{6.1}\,\%), because a constant-current charge does not
excite the capacity and the parameter loop drifts. The parameter update must be gated on
excitation --- a minimum SOC swing since the last anchor, or a minimum innovation-to-noise
ratio --- and the capacity carried from one flight to the next without reset. The joint EKF is
the recommendation; the dual EKF costs the same and is four times worse in SOC.

\paragraph{Rule 4: handle sensor faults with covariance adaptation or a bounded influence
function, not with a switching gain.} The $\mR$-adapting filters and the M-estimator filters
(Huber, maximum-correntropy, Student's~$t$) recover nominal accuracy under a noise step,
\SI{5}{\percent} outliers and a stuck channel, at EKF cost. The strong-tracking filter and
the SVSF, whose gains are driven by the innovation magnitude, are the two estimators most hurt
by sensor noise, and the SVSF is, with the ensemble EnKF, the only filter of the Kalman
families that diverges anywhere in the benchmark. Hard gating (rejecting an innovation beyond $k\sigma$) belongs in the fault
manager, not in the estimator.

\paragraph{Rule 5: expect a structurally wrong model at high C-rate, and design for it.}
Every converged Kalman-type filter is confidently wrong on the electrochemical plant
(\cref{ch:res-cell_spm}), and the voltage innovation does not reveal it. Two responses are
available and should be combined. In the estimator, bound the memory: the fading-memory EKF
holds \num{0.36}\,\% where the EKF holds \num{3.7}, and every fixed-gain observer holds
\num{0.6}. In the system, run a second estimate that does not share the model --- Coulomb
counting from the calibrated sensor with OCV re-anchoring at rest costs
\SI{0.2}{\micro\second} and \num{0.17}\,\% on the nominal mission --- and raise a flag when the
two disagree by more than their combined uncertainty. Better still, use a model that describes
the cell under load: a reduced-order electrochemical model in the estimator is the natural
next step for a cell that sees \SI{6}{C} pulses (\cref{ch:res-conclusions}).

\paragraph{Rule 6: at module level, average inside the filter.} The bar--delta and
information-fusion architectures halve the per-cell error of twelve independent EKFs at a
seventh and a third of the cost (\cref{ch:res-pack}), and find the weak cell. The decentralised
architecture is for parallel groups and for fault isolation, and there the per-cell filter
should be the plain (adapted) EKF. Balancing, not estimation, fixes an unbalanced module.

\paragraph{Rule 7: report the worst column, not the mean.} A certification argument is made on
the peak error during the landing hover and on the final error at the go-around decision
(\cref{ch:metrics}), not on the RMSE of a nominal flight, on which 51 of 70 estimators are
indistinguishable. The composite ranking of \cref{tab:ranking} was built from the geometric
mean precisely so that a single catastrophic scenario is visible in it, and the ``worst''
column is the one to read first. An interval observer (\cref{ch:intervalobserver}) or the
covariance of a well-tuned filter provides the bound that such an argument needs. The interval
observer's guaranteed bounds are never violated in any run of either plant, at a cost of
\SI{0.8}{\micro\second}; they are also \SIrange{25}{37}{\percent} of capacity wide, because
they are built for the full disturbance set of the benchmark, so they serve as a sanity
envelope and a divergence detector rather than as the reserve figure itself, which comes from
the filter's covariance.

\paragraph{Rule 8: build for \code{float}, keep it deterministic, and do not buy arithmetic
you cannot use.} No estimator in the library needs double precision on this problem
(\cref{ch:res-numerics}); the square-root forms are insurance for larger state vectors. Below
\SI{2}{\micro\second} per step the cost is irrelevant on any flight processor, and above it
--- the Gauss--Hermite filter, the particle and ensemble filters, MHE --- the benchmark finds no
accuracy that a cheaper estimator does not also deliver, with the single exception of MHE's
bias rejection. The particle filters are the correct tool for a multimodal posterior, which
the SOC problem does not have.

\section{A reference architecture}

\Cref{tab:res-arch} assembles the rules into one estimator stack for a series-module,
per-cell-voltage-sensing BMS of an electric aircraft. Every element of it is in the library and was
run in the benchmark; for a twelve-cell module the stack costs between \SI{15}{} and
\SI{30}{\micro\second} per \SI{10}{\hertz} step on the benchmark machine, depending on
whether capacity is identified per cell or on the module-average cell --- of the order of a
millisecond on a Cortex-M7, or one percent of one core.

\begin{table}[H]\centering\small\setlength{\tabcolsep}{4pt}
\caption{Reference estimator stack for an electric-aircraft BMS, assembled from the benchmark.
Costs are measured per step on the benchmark machine; the last column names the chapters that
document the element and the scenarios that motivate it.}
\label{tab:res-arch}
\begin{tabular}{@{}>{\raggedright\arraybackslash}p{0.17\textwidth}>{\raggedright\arraybackslash}p{0.36\textwidth}>{\raggedright\arraybackslash}p{0.12\textwidth}>{\raggedright\arraybackslash}p{0.27\textwidth}@{}}\toprule
element & implementation & cost & evidence \\ \midrule
per-cell SOC & EKF with innovation-based $\mR$ adaptation, fading memory bounded at a
  \SIrange{20}{60}{\second} memory length, SOC-variance floor, Joseph form, \code{float} build &
  \SI{0.5}{\micro\second} & \cref{ch:iaeakf,ch:fadingekf}; \code{noise\_step}, \code{outliers},
  \code{combined}, SPM plant \\
module fusion & bar--delta (one full filter on the average cell, twelve small deviation
  filters) or information fusion of twelve local filters & \SIrange{0.8}{1.9}{\micro\second} per
  module & \cref{ch:bardelta,ch:infofusion}; \code{weak\_cell}, cost \\
capacity (SOH) & joint EKF with $[Q, R_0]$ as parameters, excitation-gated parameter update,
  carried across flights & \SI{1.7}{\micro\second} & \cref{ch:jointekf}; \code{aged\_cell},
  multi-flight study, ground-charge profile \\
independent monitor & Coulomb counting on the calibrated current with OCV re-anchoring
  under light load; disagreement flag against the filter & \SI{0.2}{\micro\second} &
  \cref{ch:coulombocvhybrid}; SPM plant, \code{param\_mismatch} \\
bound for the safety case & interval observer or the adapted filter's $3\sigma$; peak and
  final error logged per flight & \SI{0.8}{\micro\second} & \cref{ch:intervalobserver};
  \cref{ch:metrics} \\
sensor management & calibrated current sensor (offset and gain), innovation gating and stuck-channel
  detection in the fault manager, not in the filter & --- & \code{current\_bias},
  \code{sensor\_dropout}, pack \code{current\_bias} \\
post-flight & unscented RTS smoother over the flight log for maintenance and for re-tuning &
  \SI{2.5}{\micro\second} per sample, offline & \cref{ch:urts}; \code{combined} \\
\bottomrule\end{tabular}
\end{table}

\section{Where a different choice is right}

The reference stack is not the only defensible one, and the benchmark identifies the cases in
which another element wins.

If the current sensor cannot be calibrated to better than a few percent of \SI{1}{C}, MHE in its
real-time-iteration form is the only estimator that holds the SOC to \num{0.1}\,\% under a
\SI{5}{\percent} offset, at \SI{17}{\micro\second} per step and \SI{27}{\kilo\byte}; the
disturbance observer is the cheap alternative at \num{0.18}\,\% and \SI{0.6}{\micro\second}.

If the cell sees high C-rate pulses and the estimator must run an equivalent circuit, the
fixed-gain observers --- the adaptive observer, the UIO, the high-gain observer, the certified
Luenberger observer --- are as good as the fading-memory EKF on the electrochemical plant and
have a stability proof; their price is a \SI{5}{\minute} convergence from a wrong initial SOC
and no covariance.

If the processor is a Cortex-M4 class cell-monitor slave with a few kilobytes of RAM per cell,
the reduced-order EKF (\SI{0.33}{\micro\second}, \SI{4.3}{\kilo\byte} including the OCV table,
\SI{0.26}{} geometric mean, no failure column) and the super-twisting observer
(\SI{0.14}{\micro\second}) are the Pareto-optimal choices.

If the posterior is genuinely multimodal --- an unknown initial SOC on a cell with an OCV
plateau, or a cell whose chemistry has a two-phase region --- the regularised or auxiliary
particle filter is the correct tool and the Gaussian-sum filter the cheaper approximation; on
the unimodal problem of this benchmark they buy nothing.

If a learned correction is wanted, the NN-EKF structure (a residual network inside an EKF) is
the one that preserves the physics and the covariance; its benchmark result shows both its
value (\num{0.10}\,\% on \code{preisach}, the best of the field) and its danger (a confident
bias on a cell that differs from the training population).

\section{What the benchmark does not settle}

Three questions remain open. The absolute numbers depend on the representative parameters of
the benchmark cell; the ranking of the estimators by their reaction to each fault is the
transferable content, and it should be repeated on a characterised cell before any of the
above is frozen. The electrochemical plant was run with an equivalent-circuit estimator only;
the obvious experiment --- the same filters on a reduced electrochemical model --- is the
first item of future work. And the timings are projected, not measured, on a flight processor.
